\documentclass[UTF8,11pt,a4paper]{ctexart}
\usepackage[margin=2.1cm,headheight=16pt]{geometry}
\usepackage{amsmath,amssymb,booktabs,array,tabularx,ragged2e}
\usepackage{enumitem,xcolor,fancyhdr,listings,tcolorbox}
\tcbuselibrary{breakable}
\usepackage{hyperref}
\usepackage{bookmark}

\newcommand{\DocVersion}{3.0}
\newcommand{\EngineVersion}{1.1}
\newcommand{\DocDate}{2026-11-23}
\definecolor{ink}{HTML}{1F2937}
\definecolor{accent}{HTML}{92400E}
\definecolor{muted}{HTML}{4B5563}
\definecolor{linkink}{HTML}{174678}
\definecolor{codebg}{HTML}{F5F7FA}
\hypersetup{unicode,colorlinks=true,linkcolor=linkink,urlcolor=linkink,
  pdftitle={SOE104 CW2 决策模型 v\DocVersion},
  pdfauthor={Group A--E},pdfsubject={Engine and Marking v\EngineVersion}}
\ctexset{section={format={\Large\bfseries\color{accent}}},
  subsection={format={\large\bfseries\color{ink}}}}
\pagestyle{fancy}
\fancyhf{}
\fancyhead[L]{\small\color{muted}SOE104 CW2 · Oregon Trail}
\fancyhead[R]{\small\color{muted}决策模型 v\DocVersion\ · Engine v\EngineVersion}
\fancyfoot[C]{\small\color{muted}\thepage}
\renewcommand{\headrulewidth}{0.4pt}
\setlength{\parskip}{0.35em}
\setlength{\emergencystretch}{1.5em}
\renewcommand{\arraystretch}{1.23}
\setlist{itemsep=2pt,topsep=3pt,parsep=0pt,leftmargin=1.65em}
\newcolumntype{Y}{>{\RaggedRight\arraybackslash}X}
\newcolumntype{L}[1]{>{\RaggedRight\arraybackslash}p{#1}}
\newtcolorbox{keybox}[1]{breakable,title={#1},colback=orange!4,
  colframe=accent,boxrule=0.55pt,arc=2pt,fonttitle=\bfseries}
\newtcolorbox{warnbox}[1]{breakable,title={#1},colback=red!3,
  colframe=red!45!black,boxrule=0.55pt,arc=2pt,fonttitle=\bfseries}
\lstset{basicstyle=\small\ttfamily,backgroundcolor=\color{codebg},
  numbers=left,numberstyle=\scriptsize\color{muted},numbersep=7pt,
  xleftmargin=2em,frame=l,rulecolor=\color{linkink!45},
  breaklines=true,columns=fullflexible,keepspaces=true,showstringspaces=false,
  aboveskip=0.65em,belowskip=0.65em}
\newcommand{\src}[1]{\textup{\small\textcolor{muted}{[#1]}}}
\newcommand{\api}[1]{\texttt{#1}}
\newcommand{\hp}[1]{K_{#1}}
\newcommand{\fl}[1]{L_F(#1)}
\newcommand{\ml}[1]{L_M(#1)}
\newcommand{\fg}[1]{G_F(#1)}
\newcommand{\ox}{V_O}

\title{\vspace{-1.3cm}\textbf{CW2 决策模型}\\[3pt]
  \large 配速、事件与成员行动的推导和实现}
\author{Group A--E}
\date{\DocDate\quad 文档 v\DocVersion\quad Engine / Marking v\EngineVersion}

\begin{document}
\maketitle
\thispagestyle{fancy}

\begin{keybox}{先看这一页}
每回合先判断候选方案的死亡、断粮和超时风险，再用英里当量比较能执行的方案。
旅途中使用可校准的资源价格。越接近终点，健康与实物库存越适合按终局分数评估。
同一套公式可以随队伍人数、医生状态和库存变化重算。
\end{keybox}

\noindent 实现时先读第 \ref{sec:contract} 节，查事件时直接看第 \ref{sec:events} 节。
第 \ref{sec:tests} 节提供边界用例和回放校准顺序。

\smallskip
\noindent\textbf{证据标记}
\src{TU} 表示 1.1 更新公告，\src{ER} 表示未被更新替换的引擎规则，
\src{API} 表示公开接口，\src{DER} 表示本文推导，
\src{EST} 表示模型参数或策略选择，\src{PENDING} 表示需通过回放或官方说明确认的细节。
资料位置见末节。文档和引擎分别编号。

\setcounter{tocdepth}{1}
\tableofcontents

\clearpage
\section{终局计分}\label{sec:objective}

记 $D$ 为终局里程，$N$ 为\textbf{初始人数}，$h_i(t)$ 为初始成员 $i$ 在时点 $t$
的健康值，死亡后记作 $0$。总健康为 $H_t=\sum_{i=1}^N h_i(t)$。
食物、牛、弹药和零件分别记为 $F,O,A,P$，终局实物总量 $U=F+O+A+P$。
现金 $M$ 只具备途中购买力。 \src{TU}

\[
\begin{aligned}
\mathrm{Game} &=40(D/2000)+30(H_T/100N)+30g(U),\\
S&=0.20\,\mathrm{OOP}+0.15\,\mathrm{Naming}+0.15\,\mathrm{Efficiency},\\
\mathrm{FINAL}&=S+0.01D+(0.15/N)H_T+15g(U).
\end{aligned}
\]

公告给出的货物项为 $g(U)=U/1000$，平台在 $U\ge 1000$ 时是否另外封顶仍待确认。
对固定实现的单次正常运行，下式给出游戏项对总分的账面增量。
代码效率评价、三种子汇总及编译、崩溃、超时等评分门槛另按题面处理。 \src{TU} \src{ER}
\[
\begin{gathered}
 \Delta H_t=H_t-H_{t-1},\qquad H_T=H_0+\sum_{t=1}^{T}\Delta H_t,\\
 \Delta(0.5\,\mathrm{Game})_t=
 0.01\Delta D_t+\frac{0.15}{N}\Delta H_t
       +15\bigl[g(U_t)-g(U_{t-1})\bigr].
\end{gathered}
\]

回血对应正的 $\Delta H_t$，受伤对应负值。满血开局时 $H_0=100N$。
当两个库存端点都低于 $1000$ 时，末项可写成 $0.015\Delta U_t$。
例如四人队增加 8 英里、总健康下降 16 HP、库存不变时，
账面增量为 $0.08-0.60=-0.52$ 分。同样的血量损失如果之后补回，
终局健康项会相应回升。 \src{DER}

\begin{center}\small
\begin{tabularx}{\linewidth}{@{}L{2.8cm}L{3.1cm}Y@{}}
\toprule
终局增量 & 对 FINAL 的贡献 & 按 $0.01$ 分每英里换算 \\
\midrule
1 英里 & $0.01$ 分 & 1 英里 \\
1 点总健康 & $0.15/N$ 分 & $15/N$ 英里。四人队为 3.75 英里。 \\
1 单位实物 & $0.015$ 分 & 在线性区间为 1.5 英里。 \\
1 元现金 & 0 分 & 到达前仍可用于采购或规避事件。 \\
\bottomrule
\end{tabularx}
\end{center}

\noindent 一名成员死亡使总健康从其当下血量降到零。四人队相比
“该成员满血到达”的终局上界少 3.75 分，也会失去其后续行动和职业能力。

\section{每回合如何选}\label{sec:policy}

\begin{enumerate}
  \item 建立合法候选，按事件时点的库存和健康裁剪收益与损失。
  \item 逐个结果分支检查全队覆灭、成员死亡、下次供餐和预计到达。
  在可规避的情况下优先排除严重风险。所有方案均有风险时，
  先最小化覆灭概率，再比较预期死亡人数，接着比较下一回合断粮与到达风险。
  预测存在区间时用保守边界，保留不确定标记。
  \item 在余下候选中选一步期望成本最低的方案。成本近似相同时，优先规划余量更大者。
  临近终点或可可靠预测结束状态时，用预期终局分重估。
\end{enumerate}

标量成本使用\textbf{正数表示损失}的统一约定。
对候选 $a$ 的每个可能结果先算实际净变化，再按概率求和。
\[
C(a)=\mathbb{E}\!\left[L_D-G_D+\sum_x\lambda_x(L_x-G_x)\right].
\]
这里 $L_x,G_x$ 是各资源的非负损失与增益，距离单位为英里。
用运行期价格与短期预测得到的是一步策略，策略表现以配对回放检验。 \src{DER} \src{EST}

\begin{center}\small
\begin{tabularx}{\linewidth}{@{}L{2.9cm}L{2.4cm}Y@{}}
\toprule
途中资源 & 基准英里当量 & 解释与调整依据 \\
\midrule
食物 $\lambda_F$ & 1 / 单位 & 低于下一餐或补给不确定时，先触发供餐检查。 \\
健康 $\lambda_H$ & 2 / HP & 校准起点。按预计可恢复性调节，临近终点可参考 $15/N$。 \\
现金 $\lambda_M$ & 0.5 / 元 & 公平价 2 元买 1 食物时的参照。交易机会减少时可下调。 \\
弹药 $\lambda_A$ & 1.5 / 单位 & 公平价 3 元乘参考现金价，保留猎取下一餐所需弹药。 \\
零件 $\lambda_P$ & 15 / 单位 & 公平价 30 元乘参考现金价，有 Mechanic 时修理用途下降。 \\
失一头牛 $\ox$ & 60 / 头 & 表格用基准。实际按失牛前后速度、存粮和到达风险评估。 \\
\bottomrule
\end{tabularx}
\end{center}

\noindent 这些价格是待校准参数。供餐和可达性检查优先处理临界库存，
结算比较时选择一种口径，避免把同一单位资源的运行期价值与终局价值相加。 \src{EST}

\section{回合顺序、配速与口粮}\label{sec:turn}

\begin{keybox}{一次回合的数据时点}
\centering
\api{decide} $\longrightarrow$ 行进 $\longrightarrow$ 喂食
$\longrightarrow$ 事件 $\longrightarrow$ \api{act} $\longrightarrow$ 结算
\par\smallskip\RaggedRight
\api{decide} 读回合开头，\api{resolve} 接到喂食后的事件，
\api{act} 读事件结算后的 \api{Situation}。事件扣的食物会影响接下来的补给与喂食。
\end{keybox}

\subsection{配速与剩余回合}

回合号 $t$ 从 $1$ 到 $120$。策略阶段还可安排 $121-t$ 次行进，
事件阶段当前行进已经过去，未来配速行进至多 $120-t$ 次，
当前成员行动仍在事件之后。以下预算均受提前结算和全队生存约束。 \src{ER} \src{API}
\[
 k(O)=\min(1,O/4),\qquad
 v_b=b\,k(O),\qquad b\in\{12,20,28\}.
\]
计算时使用实数或整数交叉乘法。与引擎逐回合对照时需确认非整数里程的取整。
相对 \api{NORMAL}，\api{GRUELING} 的一步净成本近似为
\[
  q_H m\lambda_H+p_O\ox-8k(O),\qquad
  (q_H,p_O)=
  \begin{cases}
  (2,0.07), & \text{有存活 Driver},\\
  (4,0.15), & \text{其他情况}.
  \end{cases}
\]
其中 $m$ 是当时活人数。四人、四牛、无 Driver 的基准下，
损失项 $32+9=41$，增加 8 英里，净成本 33。医生光环只影响事件伤害。 \src{ER} \src{TU} \src{DER}

\subsection{逐回合重算可达性}

令 $D_{\rm now}$ 是所在阶段的已走里程，$\ell=\max(0,2000-D_{\rm now})$。
为每个候选配速或事件结果搭建一个余程计划，用已规划的配速里程与成员净里程相加，
再减去预留的绕道里程 $B$。定义预测余量为
\[
 J=\sum_{\text{未执行的行进}}v_b
   +\sum_{\text{未执行的动作阶段}} v_{\rm act}
   -B-\ell .
\]
事件自身的里程结果直接调整 $D_{\rm now}$。牛损失按发生时点更新后续行进。
喂食刚完成时，仍可给本回合剩余动作安排里程。
$v_{\rm act}$ 按可持续的角色分工估计，负责狩猎、治疗或防守的成员同回合贡献 0。
预算中可先采用保守情景，再记录更快配速或补牛情景。只有存在供餐与健康预算可支撑、
且 $J\ge0$ 的计划时，才把该候选归为预计可达。结果每回合重算。 \src{EST}

三头牛时，\api{NORMAL} 基础为 15 英里，\api{GRUELING} 为 21 英里，
活着的 Scout 用 \api{SCOUT} 还能加 6。比较失牛前后的预算，
若失一头牛会使当前可行计划需要提速、补牛或增加行动里程，就重算其成本。
若替代计划也无法满足预算，将该结果标记为到达风险。商人尚未出现时，
把未来买牛作为有条件的补救情景。 \src{DER} \src{EST}

\subsection{供餐、库存与饥饿}

\begin{center}\small
\begin{tabularx}{\linewidth}{@{}L{3.1cm}L{2.7cm}Y@{}}
\toprule
口粮 & 每人食品 $r$ & 充分供餐时的健康效果 \\
\midrule
\api{BARE} & 1 & 每人 $-1$ HP \\
\api{MEAGER} & 3 & 健康不变 \\
\api{FULL} & 6 & 每人最多 $+1$ HP，上限 100 \\
\bottomrule
\end{tabularx}
\end{center}

成员先吃，按喂食时活人数 $m'$ 判断是否满足 $F\ge rm'$。
牛随后按每头 2 单位从余粮中进食。全车食品预算为 $rm'+2O'$，
故 $\lfloor F/(rm'+2O')\rfloor$ 仅在分母为正时用作库存续航指标。
成员吃不饱时，连续饥饿伤害依次为每人 5、8、12、17、23 HP，
此后保持 23。一次吃饱重置连续饥饿计数。
饥饿当回合的 \api{HEAL} 与 \api{REST} 无效，\api{HUNT} 仍可执行。 \src{TU} \src{ER}

例如四人五牛，\api{MEAGER} 至少需 12 食物喂饱成员，全车预算为 22。
若喂食前有 14 食物，成员本轮已吃饱，余下 2 才轮到牛。
下回合的供餐检查考虑当前猎取或交易及事件扣粮，再按预计活人数重算。
配速与口粮一起枚举，避免把因饥饿而无法维持的高速计划算作可达。 \src{DER}

\subsection{开局采购}

先估计划回合数和各档口粮的用量，再检查一名 Hunter 每回合最多狩猎一次、
每次耗 2 发弹药且最多增加 10 食物，以及 500 食物的仓位。
例如四人四牛用 \api{MEAGER} 走 80 回合，理论上需 $20\times80=1600$ 食物。
开局 500 食物和 120 发弹药，即使全部由 Hunter 兑现且从不溢出，也只覆盖 $500+600=1100$。
余下 500 要由途中交易、狩猎事件、较低口粮或缩短行程补足。
对照经典篮子（5 牛、500 粮、80 弹、2 零件，余 \$300），
可测试上述 4 牛、500 粮、120 弹、2 零件、余 \$260 的篮子。
同时记录多一头备用牛每回合多耗的 2 食物，以及首头牛损失对配速的影响。 \src{ER} \src{DER} \src{EST}

\clearpage
\section{事件姿态}\label{sec:events}

\subsection{套公式之前先检查状态}

以下 $s$ 为事件严重度，$p=s/10$，$m$、$F$、$O$ 等取\textbf{事件发生时}的值或估计。
食物损失 $\fl{x}=\min(x,F)$，获得量 $\fg{x}=\min(x,500-F)$，
零件损失按库存截断为 $I_P=\min(1,P)$，牛损失在 $O=0$ 时为零。
列中 $\ml{x}$ 表示被盗的有效现金数，基准表按现金充足计算，
低库存时暂按 $\min(x,M)$ 估计。 \src{ER} \src{EST}

设 $d=0.5$（事件时有存活 Doctor）或 $1$（没有）。
原始伤害每人 $x$ HP 的分支，在健康非临界区以
$\hp{x}\approx m d x\lambda_H$ 估价。精确损失应用
$\lambda_H\sum_i\min(h_i,r_d(x))$，其中 $r_d(x)$ 是引擎实际整数伤害。
取整待确认，生存检查暂用 $\lceil dx\rceil$ 作保守上界。
如医生是否存活不确定，分别估计有无光环。
先查致死分支，再使用下表的线性近似。 \src{TU} \src{DER} \src{PENDING}

\subsection{通用成本表}

表列的是事件直接成本，单位为英里。表中未在更新公告改动的分支沿用规则说明。
无法支付渡船费或逃跑费时，两种 \api{SAFE} 都按引擎转成 \api{AVOID}。
杀伤、失牛和供餐风险按第 \ref{sec:policy} 节覆盖表中排序。

\begin{center}\small
\begin{tabularx}{\linewidth}{@{}L{2.65cm}Y Y L{2.05cm}@{}}
\toprule
事件 & SAFE & RISKY & AVOID \\
\midrule
\api{RIVER} &
  $(20+s)\lambda_M$ &
  $p(\hp{10}+\lambda_F\fl{20})+0.20\ox$ & 5 \\
\api{STORM} &
  $\hp{s}$ & $\hp{s}+0.15\ox$ & 10 \\
\api{BREAKDOWN} &
  有 Mechanic 为 0。有零件时为 $\lambda_P$，否则为 15。 &
  $p(\hp{5}+\lambda_P I_P)$ & 15 \\
\api{ILLNESS} &
  随严重度变化，完整成本待确认。 & $p\hp{s+5}$ & $\hp{s}$ \\
\api{ATTACK} &
  $(20+s)\lambda_M+\lambda_F\fl{10}$ &
  $\hp{s}$ & 15 \\
\api{HUNT\_OPP} &
  $2\lambda_A-\lambda_F\fg{20}$，需至少 2 发弹药 &
  $4\lambda_A-\lambda_F\fg{50}+p\hp{5}$，需至少 4 发弹药 & 0 \\
\api{BAD\_TRAIL} &
  $6+s+\lambda_F\fl{5}$ &
  $p(\hp{s}+0.15\ox)-5(1-p)$ & 20 \\
\api{THEFT} &
  $\lambda_F\fl{5+s}$ &
  $\tfrac12\hp{5}+\tfrac12(\lambda_F\fl{20}+\lambda_M\ml{20})$ &
  $\lambda_M\ml{15}$ \\
\api{TRADER} & 挂牌成交 & 议价成交，失败不交易 & 0 \\
\bottomrule
\end{tabularx}
\end{center}

\noindent\src{TU} \src{ER} \src{DER}
\api{RIVER} 按规则文字将 20\% 失牛视作涉水后的独立检定，
\api{BAD\_TRAIL} 的 15\% 失牛只放在 $p$ 的失败分支。
\api{BAD\_TRAIL} 的 \api{SAFE} 保留旧版的 5 食物损失，公告只列了改动的里程项。
这三个分支细节可用回放核验。医生光环可能改变姿态排序。
例如 $s=2$ 的风暴，四人队有医生时预估健康损失价值为 8，
没有医生时为 16，绕道成本为 10。 \src{PENDING} \src{DER}

\clearpage
\subsection{基准速查与关键边界}

取四名活成员、有 Doctor、库存及仓位充足，按第 \ref{sec:policy} 节基准价格比较。

\begin{center}\small
\begin{tabularx}{\linewidth}{@{}L{2.65cm}Y@{}}
\toprule
事件 & 基准首选与覆盖条件 \\
\midrule
\api{RIVER} & \api{AVOID} 为 5。若损失 5 英里会使预计到达不可行，重算可支付的渡船。 \\
\api{STORM} & $s\le2$ 取 \api{SAFE}，$s\ge3$ 取 \api{AVOID}。
相同伤害下，\api{RISKY} 额外承担失牛风险。 \\
\api{BREAKDOWN} & 有 Mechanic 取 \api{SAFE}。
无 Mechanic、有零件时，$s\le4$ 取 \api{RISKY}，$s\ge5$ 时 \api{AVOID} 与 \api{SAFE} 同价。
无零件时应重算实际损失，\api{SAFE} 等效于 \api{AVOID}。 \\
\api{ILLNESS} & 先做生存检查，再比较确定伤害和随机重伤。未公布的 \api{SAFE} 需单独校验。 \\
\api{ATTACK} & $s\le3$ 取 \api{RISKY}，$s\ge4$ 取 \api{AVOID}。 \\
\api{HUNT\_OPP} & 满足弹药、仓位与生存条件时，\api{RISKY} 的成本为 $-44+2s$，
小于 \api{SAFE} 的 $-17$。空间不足时重新比较。 \\
\api{BAD\_TRAIL} & $s=1$--$5$ 取 \api{RISKY}，$s=6$--$8$ 取 \api{SAFE}，
$s=9$ 时 \api{SAFE}/\api{AVOID} 同价，$s=10$ 取 \api{AVOID}。 \\
\api{THEFT} & $s\le2$ 取 \api{SAFE}，$s\ge3$ 取 \api{AVOID}。 \\
\api{TRADER} & 按报价、成交概率及实购量重算，见下节。 \\
\bottomrule
\end{tabularx}
\end{center}

\noindent \api{BAD\_TRAIL} 的冒险成本为 $0.4s^2+1.4s-5$。
例如 $s=6$ 时，\api{RISKY}、\api{SAFE}、\api{AVOID} 分别为 $17.8$、$17$、$20$。 \src{DER}

\subsection{\texorpdfstring{\api{ILLNESS}}{ILLNESS} 的致死区间}

比较 \api{AVOID} 与 \api{RISKY} 时，设 $a=r_d(s)$、$b=r_d(s+5)$。
健康为 $h$ 的成员在 $h\le a$ 时会因 \api{AVOID} 死亡，
在 $a<h\le b$ 时仅会被 \api{RISKY} 的失败分支杀死。
当 \api{AVOID} 也会杀死成员且 $s<10$，冒险有无伤机会，
仍须按覆灭概率和预期死亡人数比较。
$s=10$ 时冒险必受更重伤害，\api{AVOID} 弱支配 \api{RISKY}。 \src{DER}

两种姿态都无致死风险且忽略取整时，成本比为 $(s+5)/10$。
这时 $s<5$ 偏向 \api{RISKY}，$s>5$ 偏向 \api{AVOID}，$s=5$ 同价。
更新公告给出 \api{SAFE} 会随严重度变化，但缺乏足够数字进行完整比较。
把它列为未定量候选，低风险局可先按已知选项执行。
若两个已知选项均有重大死亡风险，可以把 \api{SAFE} 作为待验证的应急策略，
单独记录回放结果。 \src{TU} \src{PENDING} \src{EST}

\subsection{\texorpdfstring{\api{TRADER}}{TRADER} 的单笔交易}

从 \api{Event} 读挂牌买入价 $B_x$ 与卖出价 $S_x$。议价成功概率为
$q=\operatorname{clamp}(0.5+0.04(s-5),0.2,0.9)$。
每个候选用\textbf{该报价下能成交的数量} $n$，
买入受 50 单位、钱包及食物空间约束，卖出受库存及交易上限约束。
按挂牌价和折扣价分别计算 $n$。
实测前用 $0.7B_x$、$1.3S_x$ 估算折扣报价，预算上对购入采用向上取整的保守单价。
只有成功分支发生资源变化。 \src{ER} \src{API} \src{EST}
\[
\begin{array}{@{}ll@{\qquad}ll@{}}
\text{挂牌买} & n(B_x\lambda_M-\lambda_x) &
\text{议价买} & qn(0.7B_x\lambda_M-\lambda_x)\\
\text{挂牌卖} & n(\lambda_x-S_x\lambda_M) &
\text{议价卖} & qn(\lambda_x-1.3S_x\lambda_M).
\end{array}
\]
候选包括不交易，成本为 0。按用途设置数量候选，如刚好补到下餐预算，
或把牛补到 4 头。必要补给优先比较\textbf{失败时的后果}。
买牛或卖牛时按数量重算行进和饲养预算。 \src{DER} \src{EST}

\clearpage
\section{成员行动}\label{sec:actions}

每名存活成员提交一个动作。下面给出发动条件满足时的直接成本，
并以动作发生时的健康、粮食空间以及剩余进度计算可兑现的收益。
$h_{\rm self}$ 是行动者健康，$h_{\min}$ 是当时最伤的活成员健康。
濒死的行动者使用 \api{MOVE} 仍要通过死亡风险检查。 \src{TU} \src{ER}

\begin{center}\small
\begin{tabularx}{\linewidth}{@{}L{2.75cm}L{5.45cm}Y@{}}
\toprule
动作 & 成本（英里当量） & 发动条件与截断 \\
\midrule
\api{HUNT} & $2\lambda_A-\lambda_F\fg{y}$ &
$y=10$（Hunter）或 7。至少有 2 发弹药。 \\
\api{HEAL} & $3\lambda_F-\lambda_H\min(5,100-h_{\min})$ &
Doctor 治疗最伤成员。需要 3 食物且本回合吃饱。 \\
\api{HEAL} & $3\lambda_F-\lambda_H\min(3,100-h_{\rm self})$ &
其他成员治疗自己，前置条件同上。 \\
\api{REST} & $-\lambda_H\min(2,100-h_{\rm self})$ &
仅在本回合吃饱时有效。 \\
\api{SCOUT} & $-6$（Scout）或 $-3$ &
无直接 HP 和食物成本。 \\
\api{MOVE} & $\lambda_H-5$ &
自身损失 1 HP，增加 5 英里。 \\
\api{DEFEND} & $-\lambda_H\sum_i\min(3,100-h_i)$ &
当回合受威胁才有效，至多一次，对活成员治疗。 \\
\bottomrule
\end{tabularx}
\end{center}

四人都至少缺 3 HP 时，防守最多恢复 12 HP，基准价值为 24 英里。
医生光环只减事件伤害，防守与治疗收益直接按实际恢复量计算。
食物已有 499 时，Hunter 的狩猎只能增加 1 食物，
成本为 $3-1=2$，通常应选推进动作。 \src{DER}

\begin{center}\small
\begin{tabularx}{\linewidth}{@{}L{2.8cm}Y@{}}
\toprule
角色 & 四人基线的动作分工 \src{EST} \\
\midrule
Scout & 默认 \api{SCOUT}。余程已足够时，比较保命或补给动作。 \\
Hunter & 将狩猎可兑现的食物与推进相比，预留后续喂食需要，也检查弹药。
动作前 food 容量不足或收益很低时推进。 \\
Doctor & 预算允许且队伍明显受伤时考虑 \api{HEAL}，由引擎选最伤成员。
成员接口只提供团队均值，暂用平均健康估计治疗收益，优先保留下次成员口粮。 \\
Mechanic & 作为指定防守者，在受威胁且有恢复价值时考虑 \api{DEFEND}，否则推进。
若防守收益不明，以生存和可达性状态决定。 \\
\bottomrule
\end{tabularx}
\end{center}

\noindent 普通成员可在 \api{MOVE} 与 \api{SCOUT} 间用成本和自身健康挑选。
两者在基准 $\lambda_H=2$ 下同为 $-3$，血量充足且里程紧张时可选 \api{MOVE}，
保健康时可选 \api{SCOUT}。成员的 \api{Situation} 缺少全队逐人健康与原生饥饿标志，
实测发生较多无效治疗时，再把带回合号的喂食预测作为可选输入。 \src{API} \src{EST}

\clearpage
\section{实现契约}\label{sec:contract}

当前 Java 是起步版本，下面是供实现和验收使用的契约。
\api{TeamFactory} 为一次游戏构造一组成员和一份 \api{Plans}，
由 \api{TeamStrategy} 每回合发布不可变快照，\api{TeamResolver} 只读。
快照字段应通过构造函数或集中发布更新，避免逐字段写入造成半旧半新的状态。

\begin{center}\small
\begin{tabularx}{\linewidth}{@{}L{3.05cm}Y Y@{}}
\toprule
数据类别 & 包含的信息 & 读取时的含义 \\
\midrule
回合开头快照 &
回合号、里程和目标、活人数、逐人健康与职业、五种资源。 &
从 \api{PartyState} 读取准确值。角色由存活成员的实际类型确定。 \\
本轮指令 &
配速、口粮，含当前行进的剩余回合。 &
与实际返回的 \api{Directive} 来自同一次计算。 \\
事件时预测 &
行进后里程、喂食后食物与活人健康范围、预测牛数、医生和机械师可能存活情况。 &
由实际指令推算。行进疲劳可能减员，失牛会改变后续喂食。 \\
滚动预算 &
可持续配速和动作方案、绕道余量、失牛分支可达性、可能的连续饥饿计数。 &
供事件比较使用。用连续回合的投影维护饥饿计数范围。 \\
\bottomrule
\end{tabularx}
\end{center}

事件仅提供类型、严重度、回合及市场价格。
截至事件前，弹药、零件和现金可沿用回合开头值，
食物、牛数、生命情况和里程用同回合预测。快照的回合号须与 \api{event.getTurn()} 一致。
快照缺失时，市场用 \api{Resolution.decline()}。故障优先 \api{SAFE}，
疾病采用单独配置的回退姿态，其他已知事件默认 \api{AVOID}。
疾病若尚无自己的回放数据，可暂用 \api{SAFE} 作为保命假设，
把触发原因和结果单列检查。 \src{API} \src{EST}

\subsection{策略入口}

以下是流程级伪代码。诸如 \api{projectToEvent} 等是项目内函数。
读取引擎状态的入口仅包括公开的 \api{PartyState}、\api{Event} 与 \api{Situation}。
投影在 \api{GRUELING} 有失牛风险时保留未失牛、失牛两个分支。
每个分支依次算行进疲劳、存活人数、成员及牛的喂食，
再算喂食效果与连续饥饿。确定无法验证时采用区间，供后续风险比较。

\begin{lstlisting}
decide(state):
    start = snapshot(state)
    travelTurns = 121 - state.getTurn()
    for pace in [SLOW, NORMAL, GRUELING]:
        for ration in [BARE, MEAGER, FULL]:
            directive = new Directive(pace, ration)
            beforeEvent = projectToEvent(start, directive)
            assess(directive, beforeEvent, travelTurns)
    selected = selectByRiskAndCost()
    plans.publish(start, selected.directive, selected.beforeEvent)
    return selected.directive
\end{lstlisting}

\clearpage
\subsection{事件入口}

事件后的余程评估还包含当前及后续成员行动。

\begin{lstlisting}
resolve(event):
    frame = plans.read()
    if frame is absent or frame.turn != event.getTurn():
        return fallbackFor(event)
    state = frame.beforeEvent
    futureTravelTurns = 120 - event.getTurn()
    if event.isMarket():
        choices = marketChoices(event, state)
    else:
        choices = posturesFor(event.kind(), event.severity(), state)
    evaluate(choices, state, futureTravelTurns)
    return selectByRiskAndCost().resolution
\end{lstlisting}

\subsection{成员入口}

{\RaggedRight
成员按队列顺序行动，使用每次调用传入的 \api{Situation}。
用 \api{getFoodLevel()}、\api{getAmmoLevel()}、\api{getOxenCount()} 读供给，
用 \api{getPartySize()} 和 \api{getPartyHealthAverage()} 读队伍情况。
\api{isThreatened()} 用于评估防守。读取事件详情前检查 \api{hasEvent()}。
自身健康由成员的 \api{getHealth()} 读取。当前分工可仅依赖这些公开信息。 \src{API}\par}

\begin{lstlisting}
act(situation):
    options = actionsForMyRole(situation, getHealth())
    removeUnaffordableHuntAndHeal(options, situation)
    removeSuicidalMoveIfAlternativesExist(options, getHealth())
    estimateClippedBenefits(options, situation)
    return Action.of(selectForSurvivalProgressAndCost(options))
\end{lstlisting}

\clearpage
\section{验算、回放校准与待确认项}\label{sec:tests}

\subsection{先核对边界用例}

\begin{center}\small
\begin{tabularx}{\linewidth}{@{}L{4.7cm}Y@{}}
\toprule
给定输入 & 应得结果 \\
\midrule
$t=120$ & 策略阶段有本轮行进机会，事件阶段余下行进为 0，
本回合仍可安排成员行动。 \\
四人、五牛、14 食物、MEAGER & 成员吃饱，牛最多吃剩余 2 食物。
再发生扣粮事件时，不能把本轮改记为饥饿。 \\
无 Doctor，疾病 $s=4$，HP 为 4、5、9、10 &
\api{AVOID} 伤害 4，杀死 HP 为 4 的成员。
\api{RISKY} 失败伤害 9，杀死前三人。
若无其他后果，预期死亡数分别为 $1$ 与 $0.4\times3=1.2$。 \\
仓内 499 食物、Hunter、2 发弹药 &
实际只收 1 食物，直接成本为 2 英里。 \\
四人、四牛、有 Doctor、无 Driver &
\api{GRUELING} 相对 \api{NORMAL} 基准净成本仍为 33 英里。 \\
\api{BAD\_TRAIL}，$s=9$，基准库存 &
\api{SAFE} 和 \api{AVOID} 都为 20，前者少绕行 5 英里并多耗 5 食物。 \\
同一价格下，议价成功率 0.5 &
成功节省费用与失败没有补给需要分别入账，不能只比较成交单价。 \\
\bottomrule
\end{tabularx}
\end{center}

\subsection{如何校准}

固定代码版本、初始队伍和购物篮子，在相同的三个平台种子上建立基线。
每组只改变一项参数或规则，例如 $\lambda_H$ 先试 $1.4,2.0,3.0$，
再试绕道预留或换配速条件。记录每个种子的终局里程、总健康、货物分项、
死亡数、饥饿回合数、Game、FINAL 及可见的运行成本。
比较每个种子的差值、三种子均值与最差结果，识别少数极端事件对排序的影响。
在重复试验不再改善评分或稳定性时冻结参数。 \src{EST}

\begin{warnbox}{实测与规则确认}
优先核验疾病 \api{SAFE} 的扣粮与里程、Doctor 奇数 HP 减伤的取整、
低牛数的里程取整、失牛检定及事件费用的边界。
接着验证货物超过 1000 时的结算方式、议价取整、零牛时行动里程、
接近 2000 英里时的结算时点。1.1 公告给出事件率约从 20\% 升至 56\%，
完整的事件率与严重度分布仍需校验。
回放中出现无法区分的情形时，保留对应估计与分支，等更多证据再替换。
\end{warnbox}

\subsection{资料索引}

\begin{itemize}
  \item \src{TU} \emph{Oregon Trail --- Trail Update 1.1}，第 1--4 页。
  \item \src{ER} \emph{SOE104 Oregon Trail --- Engine Rules Spec}，第 2、4--9 节。
  \item \src{API} \emph{SOE104 Oregon Trail --- API Reference}，第 1--8 节和本地公开 API JAR。
  \item \emph{SOE104-CW2-Tasksheet}，评分、提交与运行约束。
\end{itemize}

\end{document}
